% ECONOMICS_PROBLEM: {"id": "D63-144", "title": "EFX and Pareto optimality for personalized bi-valued goods", "jel_code": "D63", "source_id": "OP-0590"}
% !TeX program = xelatex
\documentclass[11pt,a4paper]{article}
\usepackage[margin=23mm]{geometry}
\usepackage{fontspec}
\setmainfont{Latin Modern Roman}
\usepackage{amsmath,amssymb,mathtools}
\usepackage{xcolor,enumitem,booktabs,longtable,array,xurl,hyperref}
\definecolor{muted}{HTML}{53636C}
\hypersetup{colorlinks=true,urlcolor=blue,linkcolor=blue}
\setlength{\parindent}{0pt}
\setlength{\parskip}{5pt}
\newcommand{\R}{\mathbb R}
\newcommand{\E}{\mathbb E}
\newcommand{\N}{\mathbb N}
\newcommand{\ind}{\mathbf 1}
\newcommand{\Var}{\operatorname{Var}}
\newcommand{\Cov}{\operatorname{Cov}}
\newcommand{\supp}{\operatorname{supp}}
\DeclareMathOperator*{\argmin}{arg\,min}
\DeclareMathOperator*{\argmax}{arg\,max}
\newcommand{\entrymeta}[3]{{\sffamily\small\color{muted}\textbf{\texttt{#1}}\quad|\quad #2\par #3\par}}
\newcommand{\Problemref}[1]{\texttt{#1}}
\newcommand{\JELref}[2]{\texttt{#2} (JEL \texttt{#1})}
\begin{document}
\section*{EFX and Pareto optimality for personalized bi-valued goods}
\textbf{Problem ID:} \texttt{D63-144}\quad \textbf{JEL:} \texttt{D63}\par
% BEGIN ECONOMICS STATEMENT
\label{problem:OP-0590}
% ECONOMICS_PROBLEM: {"local_id": "ECON-003", "title": "EFX and Pareto optimality for personalized bi-valued goods", "kind": "P", "audit_date": "2026-10-08", "source_identity_checked": true, "agenda_or_question_support_checked": true, "source_role": "see per-entry audit and locators", "current_open_certified": false}


\label{audited:ECON-003}

{\sffamily\small\color{muted}\textbf{ECON-003} | P | Explicit published open question, at source date\par}

\subsection*{Definitions and formal target}

For every agent $i$, fix two positive numbers $a_i<b_i$; each good $g$ has $v_{ig}\in\{a_i,b_i\}$. Utilities are additive. EFX requires $v_i(A_i)\geq v_i(A_j\setminus\{g\})$ for every $i,j$ and $g\in A_j$. A complete allocation $A$ is Pareto optimal (PO) when no complete allocation $B$ satisfies $v_i(B_i)\geq v_i(A_i)$ for all $i$, strictly for at least one.

\[
 \forall n,M,(a_i,b_i),(v_{ig}),\qquad
 \exists A\in\operatorname{EFX}(v)\cap\operatorname{PO}(v)?
\]
Both sets refer to integral, complete allocations; no fractional reassignment is allowed in the PO comparison.

\subsection*{Clarifications and required assumptions}

The numbers $a_i,b_i$ may differ across agents, but are constant across items for a given agent. Thus ``two types of goods'', meaning two item value vectors shared across agents, is a different restriction. The main quantifiers range over all finite $n\geq2$, all finite $M$, all positive pairs $a_i<b_i$, and all matrices respecting them. PO refers to every complete integral partition, including partitions that fail EFX.

\subsection*{Variants and changes of scope}

\textbf{Common ratio.} If $b_i/a_i=r$ for every $i$, positive row scaling converts the domain into common bi-valued values $\{1,r\}$; this is a restricted benchmark. \textbf{Zero low value.} $a_i=0$ changes the domain and the EFX/EFX$_0$ distinction. \textbf{Fractional PO.} Permit comparison with $x_{ig}\geq0$, $\sum_i x_{ig}=1$, and utilities $\sum_gv_{ig}x_{ig}$; this strengthens integral PO. \textbf{Algorithm.} Ask for polynomial time in $n,m$ and the rational input length, separately from existence.

\subsection*{Audit conclusion and source relationship}

The draft's later status citation is genuine, but [S1] is now the direct paper that poses the joint EFX--PO question. [S2] corroborates it subsequently. Existence of EFX alone does not answer the joint question. The main domain retains strictly positive personalized values.

\subsection*{References and exact source roles}

{\footnotesize\raggedright
\par\noindent\textbf{[S1]} Jiarong Jin and Biaoshuai Tao (2025). \emph{On Pareto-Optimal and Fair Allocations with Personalized Bi-Valued Utilities}. arXiv:2507.18251v1; author preprint.\par \textit{Verified locator / source role:} Introduction, personalized bi-valued domain; concluding discussion of EFX and PO. Direct source for the question, rather than only a later paper mentioning it.\par \url{https://arxiv.org/html/2507.18251v1}\par\smallskip
\par\noindent\textbf{[S2]} Wentao He and Biaoshuai Tao (2026). \emph{EFX for Additive Chores: Nonexistence, Pareto Incompatibility, and Bi-Valued Existence}. arXiv:2606.08872v2, 9 July 2026; original submission 7 June 2026. Author preprint.\par \textit{Verified locator / source role:} Section 6 (open questions); Section 2 (definitions). This version is newer than the draft\textquotesingle s v1 citation.\par \url{https://arxiv.org/html/2606.08872v2}\par\smallskip
}

\subsection*{Common conventions: Additional-source status and mathematical conventions}

\label{audited:conventions}

Of the 100 supplied ECON entries, 65 distinct problems appear here; 32 close
matches refine earlier entries, and three duplicate questions map to existing
entries. Appendix~\ref{app:audited-crosswalk} lists every destination.
The attachment's P/R/S/U distinctions and source audits are retained. Its
precise mathematical formalizations are editorial unless the individual audit
says otherwise. Candidate subsidy targets ECON-004--006 retain their U flags;
the resolved approval-core question ECON-007 updates OP-0202 above.
The following supplied conventions govern both this part and the attributed
ECON refinements elsewhere in the catalogue, subject to local restrictions.

\textbf{Parameterized questions.} A problem family lists inputs that must be fixed before an instance is evaluated: population, horizon, policies, information, data law, model class and welfare weights. These are required choices, not quantities the citations are assumed to specify. Undefined applications should not be treated as identified estimands.

\textbf{Indices and integrability.} Sets of agents, firms and regions are finite unless a population measure is explicitly used. \(i,j\) index units, \(r\) regions, \(t\) dates and \(h\) horizons. \(\mathbb E\) and \(\Pr\) refer to the declared population and law; \(\mathbf1\{A\}\) is an indicator. Every displayed expectation and discounted sum needs existence in the stated domain. Infinite horizons use a discount in \((0,1)\) and a convergence condition; finite horizons may use discount one.

\textbf{Optimization and existence.} A displayed \(\max\), \(\min\), or \(\arg\max\) is conditional on attainment. For finite feasible menus this is immediate; general spaces need continuity and compactness/coercivity or another existence argument. Without attainment, the target is the supremum/infimum and arbitrarily near-optimal feasible choices. \(\inf\varnothing=+\infty\). Empty feasibility and an unidentified target are different issues.

\textbf{Potential outcomes.} \(Y_i(z)\) is the outcome under a specified intervention, not the observed conditional mean among people choosing \(z\). In binary comparisons \(\Delta Y_i=Y_i(1)-Y_i(0)\). Specify assignment, treatment versions, compliance, information and observation horizon. Interference is allowed under a full policy vector; compressed exposure notation needs a justified mapping. No interference, consistency and overlap are substantive assumptions, not automatic consequences of notation.

\textbf{Populations and observation.} Fix baseline eligibility and population weights. Conditioning on post-policy employment, survival, relocation or program uptake can change the population being compared. Death is not ordinary loss to follow-up; outcomes truncated by death need a composite convention, joint survival target, or explicitly defined survivor stratum. Fertility-changing interventions need a population functional rather than an unexplained fixed descendant cohort. Measurement and missingness models must be supplied.

\textbf{Identification.} A structural model \(m\in\mathcal M\) implies observable law \(P_m\) and target \(\theta(m)\). Identification means
\[
 P_{m_1}=P_{m_2}\ \Longrightarrow\ \theta(m_1)=\theta(m_2)
 \quad\text{for all }m_1,m_2\in\mathcal M .
\]
Without identification, the target is
\[
 \Theta(P;\mathcal M)=\{\theta(m):m\in\mathcal M,\ P_m=P\}.
\]
Writing \(\theta(P)\) before establishing identification can hide incompatible latent models. Confidence sets additionally need a sampling law, confidence level, asymptotic sequence and uniformity class. Point identification, coverage and informative precision are separate properties.

\textbf{Mediation.} Total effects of randomized assignment are distinct from cross-world natural indirect effects. Belief/effort-channel effects require a meaningful mediator intervention and additional measurement, exclusion or mediation assumptions. Randomization of the initial policy alone does not supply those assumptions. Report bounds or interventional alternatives when the stronger target is unsupported.

\textbf{Equilibrium.} A counterfactual \(W(z)\), \(p(z)\), or \(\mu_t^z\) requires individual optimization, feasibility, government/resource budgets, market clearing and state-transition laws. State equilibrium existence and selection, or report ranges over compatible equilibria. Initial-state integration includes subsequent endogenous transitions; current-state integration must use the policy-dependent distribution. Reduced-form invariance after policy is not assumed.

\textbf{Welfare accounts.} Scores, utility, health, dollars and ecological quantities require explicit conversion before addition. Fees, taxes, subsidies, wages and extortion payments are transfers between parties; distributional weights can make them welfare relevant, but their face value is not automatically a resource loss. State ownership, geographic boundaries, foreign-income weights and fiscal finance. Do not add profits to household welfare already including dividends, or count land capitalization and the underlying benefit twice. A benefit-cost expression must distinguish gross benefits from net welfare and subtract every real cost exactly once.

\textbf{Derivatives and risk.} Log derivatives require positive arguments; local responses require admissible perturbations and specified equilibrium branches. Differentiation under expectations needs regularity/domination, and boundaries may allow only one-sided derivatives. Risk uses a specified law; ambiguity uses a declared nonempty law set on a common history space. Adaptive policies must be nonanticipative. A robust criterion is an editorial choice unless attributed explicitly.

\textbf{Scope overlaps.} Allocation, collective choice, contracts, household bargaining and coordinated policy overlap economics with optimization and game theory. They are retained because they appeared in the original catalogue; no claim of a strictly game-theory-free selection is made.
% END ECONOMICS STATEMENT
\end{document}
